\chapter{Himpunan, Relasi, Fungsi, dan Kardinalitas}

Sebuah penguji keamanan menerima spesifikasi model: pengklasifikasi dengan tiga puluh
fitur biner, masing-masing hanya bernilai nol atau satu. Untuk menguji menyeluruh, ia
harus mencoba setiap masukan yang mungkin. Setiap fitur menggandakan jumlah
kemungkinan, sehingga jawabannya $2^{30} = 1.073.741.824$. Pada kecepatan satu juta
masukan per detik, pengujian itu memakan $1073{,}74$ detik, sekitar $17{,}9$ menit.
Tambahkan sepuluh fitur lagi: jumlahnya menjadi $2^{40} = 1.099.511.627.776$, dan
mesin yang seribu kali lebih cepat memerlukan $1099{,}51$ detik, sekitar $18{,}3$
menit. Kita menggandakan masalah seribu kali dan mempercepat mesin seribu kali, lalu
kembali ke titik awal.

Kesulitan yang sama muncul lebih tajam begitu fitur boleh bernilai sebarang bilangan
real. Jumlah masukan berhenti bersifat "banyak" dan mulai bersifat "tak terhingga",
dan pertanyaannya berubah: apakah dua himpunan tak terhingga masih dapat dibandingkan
ukurannya? Naluri mengatakan bahwa bilangan rasional lebih banyak daripada bilangan
asli, sebab di antara dua bilangan asli selalu ada tak berhingga bilangan rasional.
Naluri itu keliru, dan bab ini menunjukkan mengapa.

Bab sebelumnya memberi kita bahasa: kuantor, implikasi, pola bukti. Yang belum ada
adalah objek yang dibicarakan. Bab ini memasang perbendaharaan dasarnya, yaitu
himpunan dan operasinya, relasi, fungsi, dan satu cara membandingkan ukuran himpunan
yang tidak bergantung pada kemampuan menghitung. Perbendaharaan itu dipakai di
seluruh buku, dari subruang di Bab 6 sampai pembuktian formal di Bab 28.

\begin{untukanak}
Bayangkan dua keranjang kelereng, satu merah dan satu biru. Kita ingin tahu keranjang
mana yang lebih banyak, tetapi tidak boleh menghitung. Ambil satu kelereng merah,
pasangkan dengan satu kelereng biru, sisihkan keduanya, dan ulangi sampai salah satu
keranjang kosong: bila keduanya kosong bersamaan isinya sama banyak, bila biru bersisa
biru lebih banyak. Cara ini tidak pernah menyebut angka; ia hanya memasangkan.

Sekarang misalkan keranjangnya tak berujung, sehingga kita tidak akan pernah selesai
memindahkan kelereng. Apakah memasangkan masih bisa dipakai? Bisa, dan itulah kejutan
bab ini: dengan aturan pemasangan yang tepat, seluruh bilangan cacah dapat dipasangkan
habis dengan seluruh bilangan genap, tanpa sisa di kedua pihak, dan pasangan bilangan
dapat dipasangkan dengan satu bilangan tunggal. Kita akan tahu mengapa memasangkan
lebih kuat daripada menghitung, dan mengapa ada himpunan tak berujung yang tetap tidak
dapat dipasangkan dengan bilangan cacah, sekeras apa pun kita mencoba.
\end{untukanak}

\section{Operasi pada himpunan dan hukum De Morgan}
\tingkat{dasar}

Satu himpunan belum berguna sampai kita dapat menyusun himpunan baru darinya. Empat
cara penyusunan didefinisikan dengan pola sama: sebut syarat keanggotaan, lalu himpunan
baru adalah kumpulan semua objek yang memenuhinya. Misalkan $X$ himpunan semesta dan
$A, B \subseteq X$; dari Bab 1 kita ingat bahwa $A \subseteq B$ berarti setiap anggota
$A$ juga anggota $B$, dan kesamaan himpunan dibuktikan dengan dua inklusi.

\begin{definisi}[Gabungan, irisan, selisih, komplemen]
Gabungan $A \cup B$ adalah himpunan semua $x \in X$ yang menjadi anggota $A$ atau
anggota $B$. Irisan $A \cap B$ adalah himpunan semua $x \in X$ yang menjadi anggota
$A$ dan sekaligus anggota $B$. Selisih $A \setminus B$ adalah himpunan semua
$x \in X$ yang menjadi anggota $A$ tetapi bukan anggota $B$. Komplemen
$A^{c} = X \setminus A$ adalah himpunan semua $x \in X$ yang bukan anggota $A$.
\end{definisi}

Kata "atau" pada gabungan bersifat inklusif: $x$ boleh berada di kedua himpunan, dan
itulah syarat agar $A \cap B \subseteq A \cup B$ selalu benar, seperti dicatat Hoffman
dan Kunze di Apendiks. Himpunan kosong $\varnothing$ diperlukan supaya irisan dua
himpunan tak bersekutu tetap berupa himpunan.

\begin{proposisi}[Hukum De Morgan]\index{hukum De Morgan}
Untuk sebarang $A, B \subseteq X$ berlaku
\[
  X \setminus (A \cup B) = (X \setminus A) \cap (X \setminus B),
  \qquad
  X \setminus (A \cap B) = (X \setminus A) \cup (X \setminus B).
\]
\end{proposisi}
\par\smallskip\noindent\textit{Bukti.} Ambil $x \in X$ sebarang. Menurut definisi
komplemen dan gabungan,
\[
  x \in X \setminus (A \cup B)
  \iff x \in X \ \text{dan}\ \text{tidak}\ (x \in A \ \text{atau}\ x \in B).
\]
Menurut hukum De Morgan untuk pernyataan dari Bab 1, ingkaran "salah satu" adalah
"keduanya salah", sehingga
\[
  x \in X \setminus (A \cup B)
  \iff x \in X \ \text{dan}\ x \notin A \ \text{dan}\ x \notin B
  \iff x \in (X \setminus A) \cap (X \setminus B).
\]
Karena $x$ sebarang, kedua himpunan sama. Kesamaan kedua diperoleh dengan mengganti
setiap "dan" menjadi "atau" pada rantai yang sama.
\hfill$\blacksquare${\parfillskip=0pt\par}

Hukum itu berlaku untuk dua himpunan; yang kita perlukan nanti adalah versi untuk
sebanyak apa pun himpunan.

\begin{definisi}[Keluarga terindeks]
Misalkan $I$ himpunan indeks dan untuk setiap $i \in I$ diberikan $A_i \subseteq X$.
\begin{align*}
  \bigcup_{i \in I} A_i &= \{\, x \in X : x \in A_i \ \text{untuk suatu}\ i \in I \,\}, \\
  \bigcap_{i \in I} A_i &= \{\, x \in X : x \in A_i \ \text{untuk setiap}\ i \in I \,\}.
\end{align*}
\end{definisi}

Versi umumnya berbunyi $X \setminus \bigcup_{i} A_i = \bigcap_{i} (X \setminus A_i)$
dan $X \setminus \bigcap_{i} A_i = \bigcup_{i} (X \setminus A_i)$. Buktinya tidak
menambah gagasan baru: hanya "untuk suatu" berubah menjadi "untuk setiap" ketika
ingkaran ditarik ke dalam, persis seperti pada Bab 1.

Angka membuat definisi itu dapat diperiksa. Untuk $X = \{1,\dots,10\}$,
$A = \{1,\dots,5\}$, dan $B = \{4,\dots,8\}$, Gambar~\ref{fig:venn} menampilkan
ketiga operasi atas himpunan yang sama: gabungan beranggota $3 + 2 + 3 = 8$, irisan
beranggota $2$, dan komplemen $B$ beranggota $3 + 2 = 5$.

\begin{figure}[htbp]
\centering
\begin{tikzpicture}[scale=0.82,every node/.style={font=\footnotesize}]
  \begin{scope}
    \begin{scope}
      \clip (-0.5,0) circle (0.85);
      \fill[layerteal!30] (-1.55,-1.25) rectangle (1.55,1.25);
    \end{scope}
    \begin{scope}
      \clip (0.5,0) circle (0.85);
      \fill[layerteal!30] (-1.55,-1.25) rectangle (1.55,1.25);
    \end{scope}
    \draw[ink,thick] (-0.5,0) circle (0.85);
    \draw[ink,thick] (0.5,0) circle (0.85);
    \draw[ink,thick] (-1.55,-1.25) rectangle (1.55,1.25);
    \node at (-1.0,0) {3};
    \node at (0,0) {2};
    \node at (1.0,0) {3};
    \node at (0,-1.0) {2};
    \node[above] at (-0.5,0.86) {$A$};
    \node[above] at (0.5,0.86) {$B$};
    \node[below] at (0,-1.25) {$A \cup B$: 8};
  \end{scope}
  \begin{scope}[xshift=3.9cm]
    \begin{scope}
      \clip (-0.5,0) circle (0.85);
      \clip (0.5,0) circle (0.85);
      \fill[accent!35] (-1.55,-1.25) rectangle (1.55,1.25);
    \end{scope}
    \draw[ink,thick] (-0.5,0) circle (0.85);
    \draw[ink,thick] (0.5,0) circle (0.85);
    \draw[ink,thick] (-1.55,-1.25) rectangle (1.55,1.25);
    \node at (-1.0,0) {3};
    \node at (0,0) {2};
    \node at (1.0,0) {3};
    \node at (0,-1.0) {2};
    \node[above] at (-0.5,0.86) {$A$};
    \node[above] at (0.5,0.86) {$B$};
    \node[below] at (0,-1.25) {$A \cap B$: 2};
  \end{scope}
  \begin{scope}[xshift=7.8cm]
    \fill[layergrey!22] (-1.55,-1.25) rectangle (1.55,1.25);
    \fill[white] (0.5,0) circle (0.85);
    \draw[ink,thick] (-0.5,0) circle (0.85);
    \draw[ink,thick] (0.5,0) circle (0.85);
    \draw[ink,thick] (-1.55,-1.25) rectangle (1.55,1.25);
    \node at (-1.0,0) {3};
    \node[ink!40] at (0,0) {2};
    \node[ink!40] at (1.0,0) {3};
    \node at (0,-1.0) {2};
    \node[above] at (-0.5,0.86) {$A$};
    \node[above] at (0.5,0.86) {$B$};
    \node[below] at (0,-1.25) {$B^{c}$: 5};
  \end{scope}
\end{tikzpicture}
\caption{Tiga diagram Venn atas himpunan yang sama, dengan banyaknya anggota pada
tiap daerah. Kiri: $A \cup B$ beranggota $3 + 2 + 3 = 8$; tengah: $A \cap B$
beranggota $2$; kanan: $B^{c}$ beranggota $3 + 2 = 5$, yaitu daerah berarsir,
sedangkan anggota di dalam $B$ ditulis pudar.}
\label{fig:venn}
\end{figure}

\begin{tangan}
Bahan: dua puluh kartu nama bekas, dua potong tali sepanjang dua meter, lantai lapang.

Bentuk dua lingkaran tali yang bertumpang tindih kira-kira seperempat lebarnya. Tulis
nama lima teman penyuka kopi dan lima teman penyuka teh. Letakkan kartu kopi di
lingkaran pertama, kartu teh di lingkaran kedua, lalu pindahkan kartu nama teman yang
suka keduanya ke daerah tumpang tindih: itulah irisannya. Sebut satu nama yang tidak
suka keduanya dan letakkan kartunya di luar kedua lingkaran.

Hitung tiga angka: kartu di dalam kedua lingkaran, kartu di tumpang tindih, dan kartu
di luar lingkaran kedua. Yang harus terlihat: jumlah kartu di gabungan bukan jumlah
kartu kopi ditambah jumlah kartu teh, sebab kartu di tumpang tindih terhitung dua
kali, dan kelebihan itu tepat sebesar irisannya.
\end{tangan}

Himpunan yang paling sering kita bentuk dari $A$ adalah himpunan semua himpunan
bagiannya, dan ukurannya tumbuh jauh lebih cepat daripada $A$ sendiri.

\begin{definisi}[Himpunan kuasa]\index{himpunan kuasa}
Himpunan kuasa $\mathcal{P}(A)$ adalah himpunan semua himpunan bagian $A$, yaitu
$\mathcal{P}(A) = \{\, B : B \subseteq A \,\}$.
\end{definisi}

\kamus{himpunan kuasa}{power set}{Daftar semua cara memilih sebagian anggota sebuah kantong, termasuk memilih kosong dan memilih semuanya.}

Sebagai contoh, $\mathcal{P}(\varnothing) = \{\varnothing\}$ beranggota satu, dan
$\mathcal{P}(\{1,2\}) = \{\varnothing, \{1\}, \{2\}, \{1,2\}\}$ beranggota empat.

\begin{proposisi}
Jika $A$ berhingga dengan $|A| = n$, maka $|\mathcal{P}(A)| = 2^{n}$.
\end{proposisi}
\par\smallskip\noindent\textit{Bukti.} Untuk $B \subseteq A$ definisikan fungsi
karakteristik $\chi_B : A \to \{0,1\}$ oleh $\chi_B(a) = 1$ bila $a \in B$ dan
$\chi_B(a) = 0$ bila $a \notin B$. Pemetaan $B \mapsto \chi_B$ adalah bijeksi antara
$\mathcal{P}(A)$ dan himpunan semua fungsi $A \to \{0,1\}$: dari $\chi$ kita pulihkan
$B = \{a : \chi(a) = 1\}$, dan dua himpunan bagian berbeda memberi fungsi berbeda.
Jadi $|\mathcal{P}(A)|$ sama dengan banyaknya fungsi $A \to \{0,1\}$.

Fungsi itu kita hitung dengan induksi pada $n$. Untuk $n = 0$ hanya ada fungsi
kosong, dan $2^{0} = 1$. Misalkan $|A| = n+1$, pilih $a_0 \in A$, dan tulis
$A' = A \setminus \{a_0\}$. Setiap $f : A \to \{0,1\}$ ditentukan tunggal oleh batasan
$f|_{A'}$ dan nilai $f(a_0)$; hipotesis induksi memberi $2^{n}$ pilihan batasan dan $2$
pilihan nilai, seluruhnya $2^{n+1}$.
\hfill$\blacksquare${\parfillskip=0pt\par}

Bukti itu tidak menyebut kombinasi: bijeksi memindahkan hitungan, induksi memecahnya
menjadi dua pilihan bebas.

\begin{bacakeras}
Bacakan berpasangan, satu orang suara anggota dan satu orang suara himpunan bagian.
Mulai lambat, lalu percepat.

Anggota: Satu jeruk masuk keranjang. Saya di dalam.
Himpunan bagian: Satu keranjang masuk gudang. Saya di dalam.
Anggota: Saya barang.
Himpunan bagian: Saya wadah.
Anggota: Saya ditunjuk satu per satu.
Himpunan bagian: Saya dibuktikan sekaligus.
Anggota: Kalau saya masuk, saya satu tanda.
Himpunan bagian: Kalau saya masuk, saya satu bukti.
Anggota: Saya kecil dan banyak.
Himpunan bagian: Saya besar dan kadang hanya satu.

Ulangi sampai kalimat terakhir keluar tanpa berpikir. Baris yang diulang tiga kali
bersama: barang masuk kantong, kantong masuk kantong, dan keduanya memakai tanda
berbeda.
\end{bacakeras}

\section{Hasilkali kartesius dan relasi}
\tingkat{dasar}

Urutan penting. Himpunan $\{1,2\}$ sama dengan $\{2,1\}$, tetapi koordinat titik pada
bidang tidak boleh ditukar tanpa memindahkan titiknya. Karena itu kita memerlukan
objek yang mengingat urutan, dan definisinya harus dibuat dari himpunan saja.

\begin{definisi}[Pasangan terurut dan hasilkali kartesius]\index{hasilkali kartesius}
Pasangan terurut $(a,b)$ didefinisikan sebagai $\{\{a\}, \{a,b\}\}$. Hasilkali
kartesius $A \times B$ adalah himpunan semua pasangan terurut $(a,b)$ dengan $a \in A$
dan $b \in B$.
\end{definisi}

Definisi itu terlihat aneh, tetapi yang dibutuhkan hanya satu sifat, yaitu
$(a,b) = (c,d)$ jika dan hanya jika $a = c$ dan $b = d$. Pemeriksaan sifat itu
memerlukan empat kasus dan kita tinggalkan sebagai Latihan 4.

Ukuran hasilkali mudah ditebak, dan buktinya memakai pemotongan menurut koordinat
pertama.

\begin{proposisi}
Jika $A$ dan $B$ berhingga, maka $|A \times B| = |A| \cdot |B|$.
\end{proposisi}
\par\smallskip\noindent\textit{Bukti.} Untuk setiap $a \in A$, himpunan
$\{a\} \times B$ beranggota $|B|$, sebab $b \mapsto (a,b)$ bijeksi antara $B$ dan
$\{a\} \times B$. Keluarga $\{\{a\} \times B\}_{a \in A}$ terdiri atas himpunan
saling lepas, karena $(a,b) = (a',b')$ memaksa $a = a'$. Gabungannya $A \times B$,
jadi $|A \times B|$ adalah jumlah $|A|$ suku bernilai $|B|$.
\hfill$\blacksquare${\parfillskip=0pt\par}

Sekarang objek yang menjadi judul bab ini. Sebuah relasi hanyalah pernyataan yang
mungkin benar atau salah untuk setiap pasangan objek, sehingga dapat disimpan sebagai
himpunan pasangan yang pernyataannya benar.

\begin{definisi}[Relasi]\index{relasi}
Relasi dari $A$ ke $B$ adalah sebarang $R \subseteq A \times B$; bila $A = B$ kita
menyebutnya relasi pada $A$. Kita menulis $a \mathrel{R} b$ untuk $(a,b) \in R$.
Daerah asal $R$ adalah $\{\, a \in A : (a,b) \in R\ \text{untuk suatu}\ b \in B \,\}$
dan daerah hasilnya $\{\, b \in B : (a,b) \in R\ \text{untuk suatu}\ a \in A \,\}$.
\end{definisi}

Karena relasi pada $A$ adalah himpunan bagian $A \times A$, banyaknya relasi pada
himpunan $n$ anggota adalah $2^{n^{2}}$. Untuk $n = 3$ angkanya $2^{9} = 512$ dan
untuk $n = 4$ menjadi $2^{16} = 65.536$; keduanya diperiksa dengan skrip di bagian
lab.

Empat sifat menyaringnya menjadi objek berguna. Relasi $R$ pada $A$ disebut refleksif
bila $a \mathrel{R} a$ untuk setiap $a \in A$; simetris bila $a \mathrel{R} b$
berakibat $b \mathrel{R} a$; transitif bila $a \mathrel{R} b$ dan $b \mathrel{R} c$
berakibat $a \mathrel{R} c$; dan antisimetris bila $a \mathrel{R} b$ bersama
$b \mathrel{R} a$ hanya mungkin ketika $a = b$.

\begin{definisi}[Urutan parsial dan relasi ekuivalensi]\index{relasi ekuivalensi}
Relasi pada $A$ yang refleksif, antisimetris, dan transitif disebut urutan parsial.
Relasi pada $A$ yang refleksif, simetris, dan transitif disebut relasi ekuivalensi.
\end{definisi}

Relasi $\le$ pada $\mathbb{R}$ adalah urutan parsial, begitu pula relasi "membagi"
pada bilangan asli positif, dan $\subseteq$ pada $\mathcal{P}(A)$ memberi contoh yang
kelak menjadi kerangka subruang pada Bab 6. Kata "parsial" dipakai karena dua anggota
tidak harus dapat dibandingkan: pada $\mathcal{P}(\{1,2,3\})$ himpunan $\{1\}$ dan
$\{2\}$ tidak termuat satu sama lain. Urutan yang setiap dua anggotanya dapat
dibandingkan disebut urutan total.

\section{Relasi ekuivalensi dan partisi}
\tingkat{menengah}

Relasi ekuivalensi memotong himpunan menjadi kelompok yang tidak tumpang tindih.
Sebelum sampai ke teoremanya, kita perlu nama untuk kelompoknya.

\begin{definisi}[Kelas ekuivalensi]\index{kelas ekuivalensi}
Misalkan $\sim$ relasi ekuivalensi pada $A$ dan $x \in A$. Kelas ekuivalensi $x$
adalah $[x] = \{\, y \in A : y \sim x \,\}$.
\end{definisi}

Karena $x \sim x$, tidak ada kelas yang kosong. Sifat yang lebih menarik: dua kelas
tidak dapat bertumpang sedikit pun.

\begin{lemma}
Misalkan $\sim$ relasi ekuivalensi pada $A$ dan $x, y \in A$. Maka (i) $x \in [x]$;
(ii) $[x] = [y]$ jika dan hanya jika $x \sim y$; dan (iii) bila
$[x] \cap [y] \neq \varnothing$, maka $[x] = [y]$.
\end{lemma}
\par\smallskip\noindent\textit{Bukti.} Bagian (i) adalah pernyataan $x \sim x$, benar
karena refleksif. Untuk (ii), misalkan $[x] = [y]$. Menurut (i),
$y \in [y] = [x]$, jadi $y \sim x$ dan karena simetris $x \sim y$. Sebaliknya,
misalkan $x \sim y$ dan ambil $z \in [x]$, artinya $z \sim x$. Dari $z \sim x$ dan
$x \sim y$ kita peroleh $z \sim y$, jadi $z \in [y]$; ini memberi
$[x] \subseteq [y]$, dan argumen dengan peran bertukar memberi $[y] \subseteq [x]$.
Untuk (iii), misalkan $z \in [x] \cap [y]$. Maka $z \sim x$ dan $z \sim y$; dari
$z \sim x$ dan simetri diperoleh $x \sim z$, dan bersama $z \sim y$ serta
transitivitas diperoleh $x \sim y$, sehingga $[x] = [y]$ menurut (ii).
\hfill$\blacksquare${\parfillskip=0pt\par}

Akibatnya himpunan semua kelas membagi $A$ menjadi potongan yang saling lepas dan
menutupi seluruh $A$.

\begin{definisi}[Partisi]\index{partisi}
Partisi himpunan $A$ adalah keluarga $\mathcal{B}$ himpunan bagian tak kosong dari
$A$ sedemikian sehingga setiap $x \in A$ berada tepat di satu anggota $\mathcal{B}$.
\end{definisi}

\begin{teorema}
Setiap relasi ekuivalensi pada $A$ menentukan partisi $A$, setiap partisi $A$
menentukan relasi ekuivalensi pada $A$, dan kedua proses itu saling membalik.
\end{teorema}
\par\smallskip\noindent\textit{Bukti.} Misalkan $\sim$ relasi ekuivalensi dan
$\mathcal{B} = \{\, [x] : x \in A \,\}$. Menurut Lemma, setiap $[x]$ tidak kosong dan
dua anggotanya saling lepas atau identik; setiap $x$ berada di $[x]$. Jadi
$\mathcal{B}$ partisi.

Sebaliknya, misalkan $\mathcal{B}$ partisi $A$ dan tetapkan $x \sim y$ bila $x$ dan
$y$ berada di anggota $\mathcal{B}$ yang sama. Refleksif dan simetris jelas. Untuk
transitivitas, misalkan $x \sim y$ dan $y \sim z$, dan misalkan $B$ anggota
$\mathcal{B}$ yang memuat $y$; maka $x \in B$ dan $z \in B$, jadi $x \sim z$.

Kedua proses saling membalik: dari $\sim$ kita memperoleh partisi kelas, dan partisi
itu mengembalikan $\sim$ semula karena $x \sim y$ tepat ketika $x \in [y]$; dari
partisi kita memperoleh $\sim$, dan kelas $[x]$ adalah anggota $\mathcal{B}$ yang
memuat $x$.
\hfill$\blacksquare${\parfillskip=0pt\par}

Banyaknya kelas disebut indeks. Himpunan semua kelas ditulis $A/{\sim}$ dan disebut
himpunan kuosien, sedangkan $q : A \to A/{\sim}$ dengan $q(x) = [x]$ disebut pemetaan
kuosien. Hoffman dan Kunze mencatat bahwa setiap relasi ekuivalensi dapat dipandang
berasal dari sebuah fungsi: dua anggota setara tepat ketika fungsi itu memberi nilai
sama.

\begin{contoh}
Pada bilangan bulat, tetapkan $x \equiv y \pmod{n}$ bila $n$ membagi $x-y$. Relasi
ini refleksif karena $n \mid 0$, simetris karena $n \mid (x-y)$ berakibat
$n \mid (y-x)$, dan transitif karena $n \mid (x-y)$ bersama $n \mid (y-z)$ berakibat
$n \mid (x-z)$. Kelasnya himpunan bilangan bersisa sama, dan ada tepat $n$ kelas,
yaitu $[0], \dots, [n-1]$, sebab setiap bilangan bulat dapat ditulis $x = qn + r$
dengan $0 \le r \le n-1$ secara tunggal.
\end{contoh}

Contoh yang lebih dekat dengan pekerjaan kita: untuk $f : A \to B$, tetapkan
$x \sim y$ bila $f(x) = f(y)$. Kelas $[x]$ adalah himpunan masukan dengan keluaran sama;
dalam pembelajaran mesin itulah "satu berkas prediksi", dan banyaknya kelas sama dengan
banyaknya label berbeda.

\begin{untukanak}
Lemari pakaian di rumah biasanya tidak diurutkan menurut aturan resmi. Kaus putih
menumpuk di satu rak, kaus bergambar di rak lain, kaus olahraga di keranjang ketiga.
Tidak ada yang menulis aturannya, tetapi setiap kaus berada di tepat satu tempat.

Itulah yang dilakukan sebuah relasi ekuivalensi: ia menetapkan kapan dua barang
dianggap sekamar, lalu mengelompokkan seluruh barang. Hasilnya pembagian rapi, setiap
barang masuk tepat satu kelompok, dan kelompok berbeda tidak berbagi anggota.

Sekarang pertanyaan yang membuat penasaran. Bagaimana kalau satu kaus menurut satu
aturan masuk kelompok putih, tetapi menurut aturan lain masuk kelompok olahraga? Apakah
pengelompokan masih rapi? Kerapian itu tidak bergantung pada aturan yang dipilih,
melainkan pada tiga sifat: setiap barang sekamar dengan dirinya sendiri, sekamar
berlaku dua arah, dan sekamar dapat dipindahkan lewat barang ketiga.
\end{untukanak}

\section{Fungsi: injektif, surjektif, bijektif}
\tingkat{menengah}

Kita tetapkan sekarang arti kata "fungsi", karena bab ini bergantung pada perbedaan
halus antara kodomain dan daerah hasil.

\begin{definisi}[Fungsi]\index{fungsi injektif}
Fungsi dari $A$ ke $B$ adalah aturan $f$ yang memasangkan setiap $x \in A$ dengan
tepat satu $f(x) \in B$. Himpunan $A$ disebut daerah asal, $B$ disebut kodomain, dan
$f(A) = \{\, f(x) : x \in A \,\}$ disebut daerah hasil.
\end{definisi}

Kodomain dan daerah hasil berbeda, dan perbedaan itu penting. Fungsi
$f : \mathbb{R} \to \mathbb{R}$ dengan $f(x) = x^{2}$ berkodomain $\mathbb{R}$ tetapi
berdaerah hasil $[0,\infty)$. Hoffman dan Kunze memakai tiga nama bergantian untuk
objek yang sama: fungsi, pemetaan, transformasi.

Untuk $S \subseteq B$ kita definisikan prapeta
$f^{-1}(S) = \{\, x \in A : f(x) \in S \,\}$. Notasi $f^{-1}$ di sini tidak
mengandaikan $f$ punya invers: prapeta selalu terdefinisi, invers belum tentu ada.

\begin{definisi}[Injektif, surjektif, bijektif]\index{fungsi surjektif}\index{bijeksi}
Fungsi $f : A \to B$ disebut injektif bila $f(x) = f(y)$ hanya mungkin ketika
$x = y$; surjektif bila $f(A) = B$; dan bijektif bila kedua-duanya.
\end{definisi}

\kamus{bijeksi}{bijection}{Pemasangan sempurna: setiap barang di kotak pertama punya satu pasangan di kotak kedua, tanpa sisa di kedua pihak.}

Gambar~\ref{fig:bijeksi} membandingkan ketiganya pada dua kolom simpul: panel kiri
tidak menyisakan simpul di pihak mana pun, panel tengah menyisakan $d$ di kanan,
panel kanan mengirim dua simpul kiri ke satu simpul kanan.

\begin{figure}[htbp]
\centering
\begin{tikzpicture}[scale=0.95,every node/.style={font=\footnotesize},>=Stealth,
  kotak/.style={draw=ink,thick,circle,inner sep=1.2pt,minimum size=4.4mm}]
  \begin{scope}
    \foreach \i/\y in {1/1.0,2/0,3/-1.0} {
      \node[kotak] (a\i) at (0,\y) {\i};
    }
    \node[kotak] (ba) at (1.9,1.0) {$a$};
    \node[kotak] (bb) at (1.9,0) {$b$};
    \node[kotak] (bc) at (1.9,-1.0) {$c$};
    \draw[->,thick,layerjade] (a1) -- (ba);
    \draw[->,thick,layerjade] (a2) -- (bb);
    \draw[->,thick,layerjade] (a3) -- (bc);
    \node[below] at (0.95,-1.75) {bijektif};
  \end{scope}
  \begin{scope}[xshift=3.4cm]
    \foreach \i/\y in {1/1.2,2/0.4,3/-0.4} {
      \node[kotak] (c\i) at (0,\y) {\i};
    }
    \node[kotak] (da) at (1.9,1.2) {$a$};
    \node[kotak] (db) at (1.9,0.4) {$b$};
    \node[kotak] (dc) at (1.9,-0.4) {$c$};
    \node[kotak] (dd) at (1.9,-1.2) {$d$};
    \draw[->,thick,layerindigo] (c1) -- (da);
    \draw[->,thick,layerindigo] (c2) -- (db);
    \draw[->,thick,layerindigo] (c3) -- (dc);
    \node[below] at (0.95,-1.75) {injektif saja};
  \end{scope}
  \begin{scope}[xshift=6.9cm]
    \foreach \i/\y in {1/1.2,2/0,3/-1.2} {
      \node[kotak] (e\i) at (0,\y) {\i};
    }
    \node[kotak] (fa) at (1.9,0.8) {$a$};
    \node[kotak] (fb) at (1.9,-0.8) {$b$};
    \draw[->,thick,layerbrown] (e1) -- (fa);
    \draw[->,thick,layerbrown] (e2) -- (fa);
    \draw[->,thick,layerbrown] (e3) -- (fb);
    \node[below] at (0.95,-1.75) {surjektif saja};
  \end{scope}
\end{tikzpicture}
\caption{Tiga pemetaan antara dua kolom simpul. Kiri: setiap simpul terpakai tepat
sekali, jadi bijektif. Tengah: kolom kanan beranggota empat dan simpul $d$ tidak
mendapat panah masuk, jadi injektif tetapi tidak surjektif. Kanan: dua simpul kiri
menuju $a$, jadi surjektif tetapi tidak injektif.}
\label{fig:bijeksi}
\end{figure}

Kata "bijektif" hanya nama pendek untuk gabungan dua sifat sebelumnya. Kekuatannya
muncul dalam proposisi berikut: bijeksi adalah syarat tepat agar fungsi dapat dibalik.

\begin{proposisi}
$f : A \to B$ bijektif jika dan hanya jika terdapat $g : B \to A$ dengan
$g \circ f = \mathrm{id}_A$ dan $f \circ g = \mathrm{id}_B$. Bila $g$ ada, ia tunggal.
\end{proposisi}
\par\smallskip\noindent\textit{Bukti.} Misalkan $g$ memenuhi kedua kesamaan. Bila
$f(x) = f(y)$ maka $x = g(f(x)) = g(f(y)) = y$, jadi $f$ injektif. Untuk sebarang
$b \in B$ berlaku $b = f(g(b))$, jadi $f$ surjektif.

Sebaliknya, misalkan $f$ bijektif. Untuk setiap $b \in B$ ada tepat satu $x \in A$
dengan $f(x) = b$: keberadaannya dari surjektif, ketunggalannya dari injektif.
Definisikan $g(b) = x$ untuk $x$ itu, maka $g(f(x)) = x$ dan $f(g(b)) = b$.

Untuk ketunggalan, misalkan $g$ dan $h$ memenuhi kedua kesamaan. Maka
\[
  g = g \circ \mathrm{id}_B = g \circ (f \circ h) = (g \circ f) \circ h
  = \mathrm{id}_A \circ h = h,
\]
yang hanya memakai assosiativitas komposisi dan kedua kesamaan itu.
\hfill$\blacksquare${\parfillskip=0pt\par}

Komposisi berkelakuan baik terhadap ketiga sifat: bila $f$ dan $g$ injektif maka
$g \circ f$ injektif, dan serupa untuk surjektif serta bijektif. Kebalikannya tidak
berlaku: dengan $f : \{1\} \to \{1,2\}$, $f(1) = 1$, dan $g : \{1,2\} \to \{1\}$,
$g(1) = g(2) = 1$, komposisi $g \circ f$ bijektif, padahal $g$ tidak injektif dan $f$
tidak surjektif.

Pada himpunan berhingga ketiga sifat saling mengunci, dan itulah sebabnya "satu-satu"
dan "pada" baru dapat dibedakan pada himpunan tak berhingga.

\begin{lemma}
Misalkan $A$ dan $B$ berhingga, $f : A \to B$, dan $C \subseteq B$. Maka
$|f(A)| \le |A|$ dengan kesamaan jika dan hanya jika $f$ injektif, dan
$|C| \le |B|$ dengan kesamaan jika dan hanya jika $C = B$.
\end{lemma}
\par\smallskip\noindent\textit{Bukti.} Kita buktikan bagian kedua dengan induksi pada
$n = |B|$. Untuk $n = 0$ berlaku $B = \varnothing$, sehingga $C = B$. Andaikan benar
untuk himpunan $n$ anggota dan misalkan $|B| = n+1$; pilih $b \in B$ dan tulis
$B' = B \setminus \{b\}$.

Bila $b \notin C$, maka $C \subseteq B'$ dan hipotesis induksi memberi
$|C| \le |B'| = n < n+1$, dengan kesamaan mustahil. Bila $b \in C$, tulis
$C' = C \setminus \{b\} \subseteq B'$. Maka $|C| = |C'| + 1 \le |B'| + 1 = |B|$;
kesamaan berlaku tepat ketika $|C'| = |B'|$, yang menurut hipotesis induksi berarti
$C' = B'$, yaitu $C = B$.

Bagian pertama diperoleh dengan induksi yang sama pada $|A|$: tulis
$A = A' \cup \{a\}$ dengan $a \notin A'$ dan perhatikan
$f(A) = f(A') \cup \{f(a)\}$, sehingga $|f(A)| \le |A'| + 1 = |A|$. Kesamaan menuntut
$f(a) \notin f(A')$ dan $|f(A')| = |A'|$; yang terakhir berarti $f|_{A'}$ injektif
menurut hipotesis induksi, dan yang pertama berarti $f$ tidak mengulang nilai pada
$a$.
\hfill$\blacksquare${\parfillskip=0pt\par}

\begin{proposisi}
Misalkan $A$ dan $B$ berhingga dengan $|A| = |B|$, dan $f : A \to B$. Maka $f$
injektif, $f$ surjektif, dan $f$ bijektif adalah tiga pernyataan yang setara.
\end{proposisi}
\par\smallskip\noindent\textit{Bukti.} Bila $f$ injektif, Lemma memberi
$|f(A)| = |A| = |B|$; karena $f(A) \subseteq B$ dan keduanya berukuran sama, Lemma
memberi $f(A) = B$, jadi $f$ surjektif. Bila $f$ surjektif, maka $f(A) = B$ sehingga
$|f(A)| = |B| = |A|$, dan Lemma memaksa $f$ injektif. Karena bijektif berarti kedua
sifat sekaligus, kesetaraan ketiganya selesai.
\hfill$\blacksquare${\parfillskip=0pt\par}

Peringatan yang harus dibawa ke bab berikutnya: proposisi itu memakai kata
"berhingga" secara esensial. Pada $\mathbb{N}$, fungsi $f(x) = 2x$ injektif tetapi
tidak surjektif, dan fungsi yang memetakan $1$ ke $1$ serta $n \ge 2$ ke $n-1$
surjektif tetapi tidak injektif.

\section{Himpunan berhingga, terhitung, dan tak terhitung}
\tingkat{lanjut}

Sekarang kita kembali ke pertanyaan pembuka. Untuk membandingkan dua himpunan tanpa
menghitung kita memasangkan anggotanya, dan yang perlu ditetapkan hanyalah bahwa
pemasangan itu memakai fungsi yang dapat dibalik.

\begin{definisi}[Berkekuatan sama]
Himpunan $A$ dan $B$ berkekuatan sama, ditulis $A \sim B$, bila terdapat bijeksi
$A \to B$.
\end{definisi}

Relasi $\sim$ refleksif (identitas bijeksi), simetris (invers bijeksi adalah
bijeksi), dan transitif (komposisi dua bijeksi bijeksi). Jadi ia berperilaku sebagai
relasi ekuivalensi, tetapi tidak pada "himpunan semua himpunan", sebab kumpulan itu
bukan himpunan; ini dibahas pada kotak untuk peneliti.

\begin{definisi}[Berhingga, terhitung, tak terhitung]\index{himpunan terhitung}
Himpunan $A$ berhingga bila $A = \varnothing$ atau terdapat $n \in \mathbb{N}$ dan
bijeksi $A \to \{1,2,\dots,n\}$; dalam hal kedua kita tulis $|A| = n$. Himpunan $A$
terhitung bila $A \sim \mathbb{N}$. Himpunan tak berhingga yang tidak terhitung
disebut tak terhitung.
\end{definisi}

\kamus{terhitung}{countable}{Semua anggotanya bisa diberi nomor urut tanpa ada yang terlewat, seperti setiap penumpang mendapat nomor tiket yang berbeda.}

Nama "terhitung" tepat: bila $A \sim \mathbb{N}$, bijeksinya adalah daftar. Kita
dapat menyebut anggota pertama, kedua, dan seterusnya, dan setiap anggota muncul tepat
sekali. Yang mengejutkan adalah betapa banyak himpunan yang dapat didaftar begitu.
Pasangan bilangan, misalnya.

\begin{proposisi}[Pasangan Cantor]\index{pasangan Cantor}
$\mathbb{N} \times \mathbb{N}$ terhitung.
\end{proposisi}
\par\smallskip\noindent\textit{Bukti.} Definisikan
$\pi(m,n) = \frac{(m+n)(m+n+1)}{2} + n$ untuk $m, n \in \mathbb{N} =
\{0,1,2,\dots\}$. Tulis $T_s = \frac{s(s+1)}{2}$, yaitu banyaknya pasangan dengan
$m + n < s$. Untuk pasangan dengan $m + n = s$ berlaku $0 \le n \le s$, sehingga
$\pi(m,n) = T_s + n$ mengambil nilai $T_s, \dots, T_s + s$ tepat satu kali
masing-masing. Karena $T_s + s = T_{s+1} - 1$, blok nilai untuk diagonal $s$ adalah
$\{T_s, \dots, T_{s+1}-1\}$. Blok-blok itu bersambung tanpa celah dan tanpa tumpang
tindih mulai dari $T_0 = 0$. Jadi setiap bilangan cacah berada di tepat satu blok,
dan di dalam bloknya berasal dari tepat satu pasangan; karena itu $\pi$ bijektif.
\hfill$\blacksquare${\parfillskip=0pt\par}

Inversnya eksplisit: $s = \lfloor (\sqrt{8q+1}-1)/2 \rfloor$, $n = q - T_s$, dan
$m = s - n$. Akar kuadrat itulah satu-satunya tempat aritmetika mesin dapat menipu
kita, dan contohnya ada di bagian lab.

Gambar~\ref{fig:kisi} menggambar kisi itu: sel pada baris $m$ dan kolom $n$ berlabel
$\pi(m,n)$, dan warna menandai diagonal $m + n = s$, sehingga membaca menurut warna
memberi $0$ sampai $14$ tanpa celah dan tanpa pengulangan. Panel kanan memakai arah
baca yang sama pada daftar empat bilangan real, dengan baris $d$ yang membalik digit
diagonal sehingga $d$ tidak ada di daftar.

\begin{figure}[htbp]
\centering
\begin{tikzpicture}[scale=0.95,every node/.style={font=\footnotesize}]
  \foreach \s/\warna in {0/layerteal,1/layerorange,2/layerindigo,3/layerjade,4/layerplum} {
    \foreach \n in {0,...,\s} {
      \pgfmathtruncatemacro{\m}{\s-\n}
      \pgfmathtruncatemacro{\v}{int(\s*(\s+1)/2+\n)}
      \draw[fill=\warna!30,draw=ink!45] (\n*0.86,-\m*0.72) rectangle ++(0.8,0.66);
      \node at (\n*0.86+0.4,-\m*0.72+0.33) {\v};
    }
  }
  \draw[->,thick] (4.35,0.33) -- (4.95,0.33);
  \node[right] at (4.95,0.33) {$n$};
  \draw[->,thick] (-0.45,-0.05) -- (-0.45,-3.0) node[below] {$m$};
  \begin{scope}[xshift=6.6cm,yshift=0.45cm]
    \foreach \i/\j/\b in {1/1/0,1/2/1,1/3/1,1/4/0, 2/1/1,2/2/0,2/3/0,2/4/1,
                          3/1/0,3/2/0,3/3/1,3/4/0, 4/1/1,4/2/1,4/3/0,4/4/0} {
      \ifnum\i=\j
        \fill[accent!28] (\j*0.52-0.52,-\i*0.42) rectangle ++(0.5,0.38);
      \fi
      \draw[ink!45] (\j*0.52-0.52,-\i*0.42) rectangle ++(0.5,0.38);
    }
    \draw[accent,thick] (0.25,-0.23) -- (1.79,-1.49);
    \foreach \i/\j/\b in {1/1/0,1/2/1,1/3/1,1/4/0, 2/1/1,2/2/0,2/3/0,2/4/1,
                          3/1/0,3/2/0,3/3/1,3/4/0, 4/1/1,4/2/1,4/3/0,4/4/0} {
      \node at (\j*0.52-0.27,-\i*0.42+0.19) {\b};
    }
    \foreach \i in {1,...,4} {
      \node[left] at (-0.62,-\i*0.42+0.19) {$x_{\i}$};
    }
    \foreach \j/\b in {1/1,2/1,3/0,4/1} {
      \fill[layergrey!25] (\j*0.52-0.52,-2.34) rectangle ++(0.5,0.38);
      \draw[ink!45] (\j*0.52-0.52,-2.34) rectangle ++(0.5,0.38);
      \node at (\j*0.52-0.27,-2.34+0.19) {\b};
    }
    \node[left] at (-0.10,-2.15) {$d$};
    \node[below] at (0.52,-2.42) {digit dibalik};
  \end{scope}
\end{tikzpicture}
\caption{Kiri: kisi pasangan $(m,n)$ dengan $m + n \le 4$, setiap sel berlabel
$\pi(m,n)$ dan warna menandai diagonal $m + n = s$; lima diagonal pertama memuat
nilai $0$ sampai $14$. Kanan: tabel digit empat bilangan real, dengan sel diagonal
diwarnai dan baris $d$ membalik digit diagonal $0,0,1,0$ menjadi $1,1,0,1$.}
\label{fig:kisi}
\end{figure}

Karena setiap rasional positif adalah pasangan bilangan asli yang sudah
disederhanakan, himpunan rasional juga terhitung. Yang perlu dijaga hanyalah agar
tidak ada bilangan yang didaftar dua kali.

\begin{proposisi}
Himpunan $\mathbb{Q}_{>0}$ bilangan rasional positif terhitung.
\end{proposisi}
\par\smallskip\noindent\textit{Bukti.} Setiap $q \in \mathbb{Q}_{>0}$ dapat ditulis
tunggal sebagai $q = p/r$ dengan $p, r$ bilangan asli dan $\gcd(p,r) = 1$. Definisikan
tinggi $q$ sebagai $p + r$. Untuk setiap $k \ge 2$ hanya ada berhingga rasional
bertinggi $k$, sebab $p$ dan $r$ berada di antara $1$ dan $k-1$. Maka daftar berikut
mencacah $\mathbb{Q}_{>0}$ tanpa pengulangan: untuk $k = 2, 3, 4, \dots$ tuliskan
semua $p/r$ bertinggi $k$ dengan $\gcd(p,r) = 1$, urut menaik menurut $p$. Setiap
rasional positif muncul tepat sekali, yaitu pada tingkat $k = p + r$ dengan bentuk
sederhana yang tunggal, sehingga daftar itu bijeksi $\mathbb{N} \to \mathbb{Q}_{>0}$.
\hfill$\blacksquare${\parfillskip=0pt\par}

\begin{akibat}
$\mathbb{Q}$ dan $\mathbb{Z}$ terhitung.
\end{akibat}
\par\smallskip\noindent\textit{Bukti.} Daftar $0, 1, -1, 2, -2, \dots$ adalah
bijeksi $\mathbb{N} \to \mathbb{Z}$. Untuk $\mathbb{Q}$, cacah setiap rasional tak
nol sebagai $\pm p/r$ dalam bentuk sederhana menurut $(p+r)$ dengan tanda bergantian,
lalu letakkan $0$ di awal; satu anggota tambahan tidak merusak sifat terhitung.
\hfill$\blacksquare${\parfillskip=0pt\par}

Jadi $\mathbb{N}$, $\mathbb{Z}$, dan $\mathbb{Q}$ berkekuatan sama, meskipun
$\mathbb{Q}$ terasa jauh lebih rapat. Rapat bukan ukuran besar; ukurannya adalah ada
atau tidaknya daftar.

\begin{tantangan}
Hotel di tepi pantai punya kamar bernomor $1, 2, 3, \dots$ tanpa ujung, dan malam
itu setiap kamar terisi tepat satu tamu. Seorang pelancong datang tanpa pesanan.
Resepsionis berkata hotel penuh; manajer berkata masih ada kamar untuknya tanpa
memulangkan seorang tamu pun.

Tunjukkan aturan pemindahan tamu yang memenuhi janji manajer, lalu jelaskan mengapa
aturan itu bijeksi antara tamu lama ditambah satu tamu baru dan seluruh kamar.
Mengapa hotel yang sama tidak dapat menampung tamu baru sebanyak bilangan real?

Petunjuk 1: jangan mencari kamar kosong, sebab tidak ada. Ciptakan kamar kosong dengan
memindahkan penghuninya.

Petunjuk 2: geser penghuni kamar $n$ ke kamar $n+1$ untuk setiap $n$, lalu taruh
pelancong di kamar $1$; periksa bahwa setiap kamar terisi tepat satu orang dan tidak
ada tamu yang hilang.

Kalau mentok, kembali ke bagian Himpunan berhingga, terhitung, tak terhitung,
khususnya Definisi berkekuatan sama.
\end{tantangan}

Cantor menunjukkan bahwa $\mathcal{P}(A)$ selalu lebih besar daripada $A$, dengan trik
yang tidak bergantung pada ukuran $A$.

\begin{teorema}[Cantor]\index{diagonal Cantor}
Untuk setiap himpunan $A$ tidak ada fungsi surjektif $A \to \mathcal{P}(A)$.
\end{teorema}
\par\smallskip\noindent\textit{Bukti.} Misalkan $f : A \to \mathcal{P}(A)$ sebarang
dan definisikan $D = \{\, a \in A : a \notin f(a) \,\}$. Maka $D \subseteq A$, jadi
$D \in \mathcal{P}(A)$. Andaikan $D = f(a_0)$ untuk suatu $a_0 \in A$. Bila
$a_0 \in D$, maka menurut definisi $D$ berlaku $a_0 \notin f(a_0) = D$, kontradiksi.
Bila $a_0 \notin D$, maka $a_0 \notin f(a_0)$, yang justru syarat keanggotaan $D$,
sehingga $a_0 \in D$, kontradiksi lagi. Jadi $D$ tidak berada di daerah hasil $f$.
\hfill$\blacksquare${\parfillskip=0pt\par}

Himpunan $D$ dibentuk dengan menanyakan, untuk setiap $a$, apakah $a$ menjadi anggota
himpunan yang dipasangkan kepadanya, lalu membalik jawabannya; pembalikan itulah yang
membuat $D$ berbeda dari setiap nilai $f$. Pada $A = \mathbb{N}$ trik ini muncul sebagai
pembacaan diagonal tabel dengan setiap entrinya diubah, karena itu disebut argumen
diagonal.

\begin{akibat}
$\mathcal{P}(\mathbb{N})$ tak terhitung, dan $\mathbb{R}$ tak terhitung.
\end{akibat}
\par\smallskip\noindent\textit{Bukti.} Menurut Teorema Cantor dengan $A = \mathbb{N}$,
tidak ada surjeksi $\mathbb{N} \to \mathcal{P}(\mathbb{N})$, jadi
$\mathcal{P}(\mathbb{N})$ tak terhitung.

Untuk $\mathbb{R}$ cukup ditunjukkan bahwa $[0,1]$ tak terhitung. Setiap $x \in [0,1]$
memiliki ekspansi biner $x = \sum_{k \ge 1} b_k 2^{-k}$ dengan $b_k \in \{0,1\}$;
keberadaannya dibuktikan pada Bab 3 dan Bab 5. Pemetaan dari himpunan barisan biner ke
$[0,1]$ karena itu surjektif, dan himpunan barisan biner berkekuatan sama dengan
$\mathcal{P}(\mathbb{N})$ sebab barisan biner adalah fungsi $\mathbb{N} \to \{0,1\}$.
Bila $[0,1]$ terhitung, bayangan prapetanya di himpunan barisan biner akan terhitung
pula, bertentangan dengan paragraf sebelumnya.
\hfill$\blacksquare${\parfillskip=0pt\par}

Satu catatan kejujuran: ekspansi biner tidak tunggal, sebab
$0{,}0111\dots = 0{,}1000\dots$ dalam basis dua. Argumen di atas tidak terganggu karena
hanya memakai sifat surjektif, tetapi $|\mathbb{R}| = |\mathcal{P}(\mathbb{N})|$
menuntut alat berikut.

\begin{proposisi}[Cantor--Bernstein--Schroeder]
Bila terdapat fungsi injektif $A \to B$ dan injektif $B \to A$, maka $A \sim B$.
\end{proposisi}
\par\smallskip\noindent\textit{Sketsa bukti.} Misalkan $f : A \to B$ dan
$g : B \to A$ injektif. Tetapkan $A_0$ sebagai himpunan titik $A$ yang tidak punya
pendahulu di $B$, lalu $A_{k+1} = g(f(A_k))$, dan tulis $A^{+}$ untuk gabungan semua
$A_k$ dengan $k$ genap. Definisikan $h(x) = f(x)$ bila $x \in A^{+}$ dan
$h(x) = g^{-1}(x)$ bila $x \notin A^{+}$. Pada bagian pertama $h$ injektif karena $f$
injektif, pada bagian kedua karena $g^{-1}$ injektif, dan bayangan kedua bagian
terpisah menurut cara $A^{+}$ dibentuk, sehingga $h$ injektif. Untuk kelinearan
bayangan, ambil $y \in B$. Bila $g(y) \in A^{+}$, maka $g(y)$ berada di suatu $A_k$
dengan $k$ genap sehingga $g(y) = g(f(x))$ untuk $x \in A_{k-1}$ dan $h(x) = y$; bila
$g(y) \notin A^{+}$, maka $h(g(y)) = y$. Jadi $h$ surjektif.

Pembukuan kedalaman rantai itulah bagian yang halus, dan rinciannya dilewati di sini;
versi formal lengkapnya ada pada \emph{Mathematics in Lean} bagian 4.3.
\hfill$\blacksquare${\parfillskip=0pt\par}

Dengan alat itu kita boleh menyimpulkan $|\mathbb{R}| = |\mathcal{P}(\mathbb{N})|$:
injeksi dari $\mathcal{P}(\mathbb{N})$ ke $\mathbb{R}$ diberikan ekspansi biner, dan
injeksi sebaliknya oleh fungsi tangen. Kekuatan $\mathbb{R}$ biasa ditulis
$2^{\aleph_0}$, dengan $\aleph_0$ kekuatan $\mathbb{N}$.

\begin{untukanak}
Bayangkan mesin tiket di pintu masuk pameran. Setiap pengunjung mendapat satu tiket
bernomor: 1, 2, 3, dan seterusnya, tidak ada nomor terpakai dua kali, dan setiap nomor
yang mungkin akan suatu saat diberikan. Kumpulan benda yang bisa diberi tiket seperti
itu disebut terhitung.

Kelihatannya semua kumpulan tak berujung bisa diberi tiket. Tetapi tidak. Ada
kumpulan yang anggotanya begitu banyak sehingga sebanyak apa pun tiket yang kita
cetak, selalu ada anggota yang tidak kebagian. Mesin tiketnya tidak rusak; kumpulan
itu memang lebih besar daripada persediaan nomor.

Pertanyaan yang membuat penasaran: bagaimana kita bisa tahu sebuah kumpulan tidak
bisa diberi tiket, padahal belum semua cara dicoba? Jawabannya trik pembalikan Cantor.
Setiap kali seseorang mengaku sudah memberi tiket kepada semua anggota, kita susun satu
anggota baru yang sengaja berbeda dari setiap anggota bertiket; anggota itu pasti
terlewat.
\end{untukanak}

\section{Ketaksamaan, nilai mutlak, dan batas atas}
\tingkat{menengah}

Bagian terakhir memasang dua alat untuk Bab 3. Yang pertama mengukur jarak antara dua
bilangan tanpa memandang arah, yaitu nilai mutlak.

\begin{definisi}[Nilai mutlak]\index{nilai mutlak}
Untuk $x \in \mathbb{R}$, $|x| = x$ bila $x \ge 0$ dan $|x| = -x$ bila $x < 0$.
\end{definisi}

Jarak melalui titik ketiga tidak pernah lebih pendek daripada jarak langsung.

\begin{proposisi}[Ketaksamaan segitiga]
Untuk setiap $x, y \in \mathbb{R}$ berlaku $|x + y| \le |x| + |y|$ dan
$\bigl| |x| - |y| \bigr| \le |x - y|$.
\end{proposisi}
\par\smallskip\noindent\textit{Bukti.} Perhatikan $-|x| \le x \le |x|$ dan
$-|y| \le y \le |y|$. Menjumlahkan keduanya memberi
$-(|x|+|y|) \le x + y \le |x|+|y|$, yang menurut definisi nilai mutlak berarti
$|x+y| \le |x|+|y|$. Untuk yang kedua, tulis $x = (x-y) + y$ dan terapkan hasil
pertama: $|x| \le |x-y| + |y|$, jadi $|x| - |y| \le |x-y|$. Menukar peran $x$ dan $y$
memberi $|y| - |x| \le |x-y|$, sehingga $\bigl| |x| - |y| \bigr| \le |x-y|$.
\hfill$\blacksquare${\parfillskip=0pt\par}

Sekarang batas. Himpunan bilangan dapat terbatas di atas tanpa memiliki anggota
terbesar, dan justru kasus itulah yang paling penting.

\begin{definisi}[Batas atas dan supremum]\index{supremum}
Misalkan $S \subseteq \mathbb{R}$ tak kosong. Bilangan $u$ disebut batas atas $S$
bila $s \le u$ untuk setiap $s \in S$; bila ada, $S$ disebut terbatas di atas.
Bilangan $u$ disebut supremum $S$, ditulis $u = \sup S$, bila $u$ batas atas $S$ dan
$u \le v$ untuk setiap batas atas $v$ dari $S$.
\end{definisi}

\kamus{supremum}{supremum (least upper bound)}{Pagar paling rapat: batas tertinggi yang tidak dapat diturunkan lagi tanpa menyisakan anggota di atasnya.}

Supremum, bila ada, tunggal, sebab dua batas atas terkecil saling mendominasi. Yang
tidak dijamin definisi adalah keberadaannya: $\{x \in \mathbb{Q} : x^{2} < 2\}$
terbatas di atas dalam $\mathbb{Q}$ tetapi tidak punya supremum di sana.

\begin{catatan}[Aksioma kelengkapan]
Setiap himpunan tak kosong $S \subseteq \mathbb{R}$ yang terbatas di atas memiliki
supremum di $\mathbb{R}$. Aksioma ini tidak dibuktikan di sini; pembangunannya
beserta akibat-akibatnya adalah isi Bab 3, termasuk ciri epsilon yang mengubah
pernyataan tentang batas terkecil menjadi pernyataan tentang $\varepsilon$, dan
contoh himpunan yang memiliki supremum tanpa maksimum.
\end{catatan}

Pada himpunan berhingga supremum dan maksimum berimpit, dan menukar keduanya sering
lolos pada contoh kecil.

\section{Contoh hitung}
\tingkat{menengah}

Empat hitungan berikut dikerjakan tangan sampai angka, lalu diperiksa ulang dengan
skrip. Semuanya kecil, sehingga kesalahan aritmetika langsung terlihat.

\begin{resep}
Untuk dua himpunan hingga $A$ dan $B$:
\begin{enumerate}
  \item Hitung $|A|$, banyaknya anggota $A$.
  \item Hitung $|B|$ dengan cara yang sama.
  \item Hitung $|A \cap B|$, anggota yang muncul di kedua daftar.
  \item Jumlahkan lalu kurangi: $|A \cup B| = |A| + |B| - |A \cap B|$.
\end{enumerate}
Peringatan urutan: langkah 3 tidak boleh ditunda sampai sesudah langkah 4. Kalau
penjumlahan dikerjakan lebih dulu, anggota irisan terhitung dua kali dan hasilnya
kelebihan tepat $|A \cap B|$. Untuk $A = \{1,\dots,5\}$ dan $B = \{4,\dots,8\}$:
$5 + 5 - 2 = 8$, sama dengan angka pada diagram Venn sebelumnya.
\end{resep}

Pertama, himpunan kuasa. Untuk $A = \{1,2,3\}$ himpunan bagiannya berukuran
$0,1,2,3$ dengan cacah $1, 3, 3, 1$, seluruhnya $1 + 3 + 3 + 1 = 8 = 2^{3}$. Baris
$n = 4$ memberi $1+4+6+4+1 = 16$, $n = 2$ memberi $1+2+1 = 4$, dan $n = 1$ memberi
$1+1 = 2$: segitiga Pascal yang sama dengan Proposisi tentang $2^{n}$.

Kedua, pasangan Cantor. Dengan $\pi(m,n) = T_{m+n} + n$ dan $T_0 = 0$, $T_1 = 1$,
$T_2 = 3$, $T_3 = 6$: diagonal $s = 0$ memuat $(0,0) \mapsto 0$; $s = 1$ memuat
$(1,0) \mapsto 1$ dan $(0,1) \mapsto 2$; $s = 2$ memuat $(2,0) \mapsto 3$,
$(1,1) \mapsto 4$, $(0,2) \mapsto 5$; $s = 3$ memuat $(3,0) \mapsto 6$,
$(2,1) \mapsto 7$, $(1,2) \mapsto 8$, $(0,3) \mapsto 9$. Sepuluh nilai itu tepat
$0, \dots, 9$ tanpa pengulangan dan tanpa celah.

Arah sebaliknya lebih mudah salah. Ambil $q = 1000$. Kita cari $s$ dengan
$T_s \le 1000 < T_{s+1}$. Karena $T_{44} = 44 \cdot 45/2 = 990$ dan $T_{45} = 1035$,
maka $s = 44$, lalu $n = 1000 - 990 = 10$ dan $m = 44 - 10 = 34$. Pemeriksaan:
$\pi(34,10) = 44 \cdot 45/2 + 10 = 990 + 10 = 1000$, benar.

Ketiga, bilangan rasional. Menurut aturan tinggi, tingkat $k = 2$ memuat $1/1$;
$k = 3$ memuat $1/2, 2/1$; $k = 5$ memuat $1/4, 2/3, 3/2, 4/1$. Banyaknya suku tiap
tingkat adalah nilai fungsi totien Euler, $1, 2, 2, 4, 2, 6$ untuk $k = 2, \dots, 7$.
Kumulatifnya sampai $k = 7$ adalah $1+2+2+4+2+6 = 17$, jadi $6/1$ menempati posisi
ke-17. Karena kumulatif sampai $k = 9$ adalah $27$ dan tingkat $k = 10$ memuat
$1/9, 3/7, 7/3, 9/1$, bilangan $7/3$ menempati posisi ke-30.

Keempat, banyaknya relasi. Pada himpunan $3$ anggota ada $2^{9} = 512$ relasi, dan
hanya $5$ yang ekuivalensi: $1$ dengan satu kelas, $3$ dengan dua kelas, $1$ dengan tiga
kelas. Pada himpunan $4$ anggota ada $2^{16} = 65.536$ relasi dan $15$ yang ekuivalen.
Keempat angka itu muncul lagi sebagai empat baris terakhir keluaran skrip; barisan
$1, 2, 5, 15, 52$ dikenal sebagai bilangan Bell.

\kertas{Empat angka yang sudah diperiksa ulang. (i) $|\mathcal{P}(A)| = 2^{n}$ untuk
$n = 0,1,2,3,4$ memberi $1, 2, 4, 8, 16$.
(ii) $\pi(34,10) = 990 + 10 = 1000$, dan $\pi$ pada diagonal $s = 0,1,2,3$ memberi
$0,1,2,3,4,5,6,7,8,9$. (iii) Banyak rasional positif bertinggi $k$ adalah
nilai fungsi totien di $k$; kumulatif sampai $k = 9$ adalah $27$, jadi $7/3$ di
posisi ke-30. (iv) Relasi pada himpunan $3$ anggota: $512$ seluruhnya, $5$ ekuivalen.}

\section{Lab komputasi}
\tingkat{lanjut}

Skrip berikut menjalankan keempat hitungan itu dan menambahkan satu pemeriksaan yang
tidak dapat dikerjakan tangan: ia membandingkan invers pasangan Cantor yang dihitung
dengan pecahan mesin dan yang dihitung dengan akar bilangan bulat. Keduanya
seharusnya memberi jawaban sama; pada satu titik jawabannya berbeda.

\begin{lab}
\begin{verbatim}
from itertools import combinations
from math import gcd, isqrt, sqrt, floor

# (1) banyak himpunan bagian untuk n kecil
for n in range(5):
    A = range(n)
    subs = [S for k in range(n + 1) for S in combinations(A, k)]
    print("n =", n, "|P(A)| =", len(subs), "2^n =", 2 ** n)

# (2) pasangan Cantor pi(m, n) = (m+n)(m+n+1)/2 + n
def pi(m, n):
    s = m + n
    return s * (s + 1) // 2 + n

pasangan = [(0, 0), (1, 0), (0, 1), (2, 0), (1, 1),
            (0, 2), (3, 0), (2, 1), (1, 2), (0, 3)]
print("pi:", [pi(m, n) for m, n in pasangan])
keys = {pi(m, n) for m in range(61) for n in range(61)}
print("kunci unik pada 61x61 pasangan:", len(keys) == 61 * 61)

# (3) invers: float vs isqrt
s, n = 2 * 10**8, 2 * 10**8 - 1
q = s * (s + 1) // 2 + n
print("q =", q, "s benar =", s)
print("s dari float:", int(floor((sqrt(8 * q + 1) - 1) / 2)))
print("s dari isqrt:", (isqrt(8 * q + 1) - 1) // 2)

# (4) enumerasi Q positif menurut tinggi k = p + q
seq = []
for k in range(2, 10):
    row = [(p, k - p) for p in range(1, k) if gcd(p, k - p) == 1]
    print("k =", k, "banyak =", len(row), row)
    seq += row
print("suku ke-1..8:", seq[:8])
print("total sampai k = 9:", len(seq))

# (5) relasi pada himpunan n anggota: 2^(n*n) seluruhnya, Bell yang ekuivalen
def ekuivalen(R, n):
    refl = all((a, a) in R for a in range(n))
    sim = all((b, a) in R for a, b in R)
    trans = all((a, c) in R for a, b in R for c in range(n)
                if (b, c) in R)
    return refl and sim and trans

for n in range(1, 5):
    semua = 2 ** (n * n)
    hitung = 0
    for bit in range(semua):
        R = {(a, b) for a in range(n) for b in range(n)
             if bit >> (a * n + b) & 1}
        if ekuivalen(R, n):
            hitung += 1
    print("n =", n, "relasi =", semua, "ekuivalen =", hitung)
\end{verbatim}
\end{lab}

Keluaran yang tercetak, apa adanya:

\begin{verbatim}
n = 0 |P(A)| = 1 2^n = 1
n = 1 |P(A)| = 2 2^n = 2
n = 2 |P(A)| = 4 2^n = 4
n = 3 |P(A)| = 8 2^n = 8
n = 4 |P(A)| = 16 2^n = 16
pi: [0, 1, 2, 3, 4, 5, 6, 7, 8, 9]
kunci unik pada 61x61 pasangan: True
q = 20000000299999999 s benar = 200000000
s dari float: 200000001
s dari isqrt: 200000000
k = 2 banyak = 1 [(1, 1)]
k = 3 banyak = 2 [(1, 2), (2, 1)]
k = 4 banyak = 2 [(1, 3), (3, 1)]
k = 5 banyak = 4 [(1, 4), (2, 3), (3, 2), (4, 1)]
k = 6 banyak = 2 [(1, 5), (5, 1)]
k = 7 banyak = 6 [(1, 6), (2, 5), (3, 4), (4, 3), (5, 2), (6, 1)]
k = 8 banyak = 4 [(1, 7), (3, 5), (5, 3), (7, 1)]
k = 9 banyak = 6 [(1, 8), (2, 7), (4, 5), (5, 4), (7, 2), (8, 1)]
suku ke-1..8: [(1, 1), (1, 2), (2, 1), (1, 3), (3, 1), (1, 4), (2, 3), (3, 2)]
total sampai k = 9: 27
n = 1 relasi = 2 ekuivalen = 1
n = 2 relasi = 16 ekuivalen = 2
n = 3 relasi = 512 ekuivalen = 5
n = 4 relasi = 65536 ekuivalen = 15
\end{verbatim}

Yang perlu dibaca dari keluaran itu: lima baris pertama mengonfirmasi $2^{n}$ untuk
$n = 0$ sampai $4$; sepuluh nilai $\pi$ berurutan tanpa celah, dan $61 \times 61 =
3721$ pasangan memberi kunci unik sebanyak $3721$, yang merupakan bukti numerik, bukan
bukti matematis, sebab bukti matematisnya argumen blok $T_s$. Empat baris terakhir
menyapu seluruh $2^{n^{2}}$ relasi untuk $n = 1$ sampai $4$ dan mencacah yang
refleksif, simetris, dan transitif: hasilnya $1, 2, 5, 15$. Baris $n = 5$ tidak ada
karena menyapu $2^{25}$ relasi terlalu lama; nilainya $52$.

Ketiga, invers dengan pecahan salah. Pada $q = 20.000.000.299.999.999$ rumus
$\lfloor (\sqrt{8q+1}-1)/2 \rfloor$ dengan pecahan mesin 64-bit memberi
$s = 200.000.001$, sedangkan nilai benarnya $200.000.000$. Kesalahannya satu, dan itu
cukup untuk mengembalikan pasangan yang sama sekali berbeda: dengan $s = 200.000.001$
nilai $n = q - T_s$ menjadi negatif. Angka $8q+1$ berorde $10^{17}$, sedangkan pecahan
64-bit hanya menyimpan sekitar $16$ angka desimal dengan tepat; \texttt{isqrt} bekerja
pada bilangan bulat dan memberi $200.000.000$. Pelajarannya umum: membalik rumus
penomoran yang eksak menuntut aritmetika bilangan bulat, seperti dibahas pada Bab 22.

\begin{untukdosen}
Empat pertemuan. Pertama: operasi himpunan dan De Morgan, tiap bukti ditulis sebagai
rantai ekuivalensi, bukan dua inklusi terpisah. Kedua: hasilkali kartesius, relasi, lalu
penyaringan menjadi urutan parsial dan relasi ekuivalensi. Ketiga: kelas, partisi, dan
fungsi, dengan Lemma kelas dibuktikan di papan. Keempat: kardinalitas sampai Teorema
Cantor, ditambah sketsa Cantor--Bernstein bila waktu memungkinkan.

Dua miskonsepsi yang paling sering muncul. Pertama, mahasiswa mengira relasi yang
simetris dan transitif otomatis refleksif; penyangkal $R = \{(1,1),(2,2),(1,2),(2,1)\}$
pada $\{1,2,3\}$ biasanya cukup, dan Latihan 2 meminta penilaian klaim itu. Kedua,
mahasiswa mengira $\mathbb{Q}$ lebih besar daripada $\mathbb{N}$ karena kerapatannya;
mintalah sepuluh suku pertama enumerasi tinggi di papan, lalu tanyakan posisi $7/3$.

Aktivitas yang berhasil: bagi papan menjadi "bilangan asli" dan "rasional positif", lalu
minta dua relawan memasangkan anggota selama dua menit. Pemasangan naif akan gagal,
lalu bersama-sama susun pemasangan menurut tinggi. Mahasiswa siap lanjut bila dapat
menjelaskan mengapa daftar itu tidak melewatkan satu pun rasional, dan mengapa daftar
serupa gagal untuk $[0,1]$.
\end{untukdosen}

\section{Ringkasan}
\tingkat{dasar}

\begin{ringkasan}
\begin{itemize}
  \item Gabungan, irisan, selisih, dan komplemen didefinisikan lewat syarat
        keanggotaan; hukum De Morgan adalah hukum logika yang dipindahkan ke
        himpunan, dan versinya berlaku untuk keluarga terindeks apa pun.
  \item $|\mathcal{P}(A)| = 2^{|A|}$ untuk $A$ berhingga, lewat bijeksi ke fungsi
        karakteristik $A \to \{0,1\}$.
  \item Relasi adalah himpunan bagian $A \times B$; pada himpunan $n$ anggota ada
        $2^{n^{2}}$ relasi, dan empat sifat menyaringnya menjadi relasi ekuivalensi
        dan urutan parsial.
  \item Relasi ekuivalensi dan partisi adalah dua wajah objek yang sama; kelasnya
        selalu tidak kosong, saling lepas, dan menutupi seluruh himpunan.
  \item Fungsi bijektif tepat ketika ia punya invers dua sisi yang tunggal, dan pada
        himpunan berhingga berukuran sama ketiga sifat saling mengunci.
  \item $\mathbb{N} \times \mathbb{N}$ dan $\mathbb{Q}$ terhitung lewat penomoran
        eksplisit, sedangkan $\mathcal{P}(\mathbb{N})$ dan $\mathbb{R}$ tak terhitung
        lewat argumen diagonal Cantor.
  \item Supremum adalah batas atas terkecil; keberadaannya dijamin aksioma
        kelengkapan, sedangkan ciri epsilon dan pembangunannya dilanjutkan di Bab 3.
\end{itemize}
\end{ringkasan}

\section{Latihan}
\tingkat{menengah}

\begin{latihan}
\begin{enumerate}
  \item Buktikan $A \cap (B \cup C) = (A \cap B) \cup (A \cap C)$ dan
        $A \setminus (B \setminus C) = (A \setminus B) \cup (A \cap C)$ dengan rantai
        ekuivalensi keanggotaan.
  \item Nilai klaim ini: benar, salah, atau belum cukup bukti. "Setiap relasi pada
        himpunan tak kosong yang simetris dan transitif pasti refleksif." Bila salah,
        berikan contoh penyangkal sekecil mungkin; bila belum cukup bukti, sebutkan
        bukti tambahan apa yang diperlukan.
  \item Buktikan bahwa bila $f : A \to B$ injektif, maka $S \mapsto f(S)$ adalah
        fungsi injektif dari $\mathcal{P}(A)$ ke $\mathcal{P}(B)$. Jelaskan di mana
        sifat injektif dipakai dan mengapa kesimpulannya gagal tanpanya.
  \item Buktikan lemma pasangan terurut Kuratowski:
        $\{\{a\},\{a,b\}\} = \{\{c\},\{c,d\}\}$ jika dan hanya jika $a = c$ dan
        $b = d$. Tangani kasus $a = b$ secara terpisah.
  \item Lengkapi bagian pertama Lemma hitungan: untuk $f : A \to B$ dengan $A$
        berhingga, $|f(A)| \le |A|$ dan kesamaan berlaku hanya bila $f$ injektif.
  \item Hitung banyaknya relasi ekuivalensi pada himpunan $5$ anggota yang memiliki
        tepat dua kelas, lalu periksa jawabannya dengan Python.
  \item Hitung $\pi(7,5)$ dengan rumus pasangan Cantor, lalu tentukan $(m,n)$ dengan
        $\pi(m,n) = 1000$. Periksa kedua jawaban dengan Python.
  \item Tentukan posisi $7/3$ dalam enumerasi rasional positif menurut tinggi, lalu
        tentukan bilangan pada posisi ke-30 dan ke-41. Periksa dengan Python.
  \item Buktikan bahwa bila $A_1, A_2, A_3, \dots$ masing-masing terhitung dan
        penomoran setiap $A_n$ diberikan eksplisit, maka $\bigcup_{n \ge 1} A_n$
        terhitung.
  \item Tunjukkan $(0,1)$ berkekuatan sama dengan $\mathbb{R}$ melalui satu bijeksi
        eksplisit, dan jelaskan mengapa fungsi yang kita pilih harus dibatasi daerah
        asalnya.
\end{enumerate}
\end{latihan}

Petunjuk untuk tiga butir tersulit. Latihan 5: pada langkah induksi tulis
$A = A' \cup \{a\}$ dengan $a \notin A'$, lalu bandingkan $f(A)$ dengan
$f(A') \cup \{f(a)\}$; kesamaan ukuran menuntut $f(a) \notin f(A')$ dan
$|f(A')| = |A'|$ sekaligus. Latihan 6: banyaknya cara memisahkan $n$ anggota menjadi
dua bagian tak kosong adalah $(2^{n}-2)/2$, sebab memilih bagian pertama atau kedua
memberi pemisahan yang sama. Latihan 9: susun penomoran seluruh $A_n$ sebagai tabel dua
kolom, baca diagonal demi diagonal seperti pada bukti bahwa
$\mathbb{N} \times \mathbb{N}$ terhitung, lalu buang suku berulang.

\begin{untukpeneliti}
Yang benar-benar terbukti di sini ada empat. Teorema Cantor tidak memakai aksioma
pilihan dan bersifat konstruktif: dari sebarang $f$ kita membangun $D$ yang tidak
berada di daerah hasilnya. Ketakterhitungan $\mathbb{R}$ bergantung pada ekspansi biner
setiap titik $[0,1]$; fakta itu memerlukan kelengkapan dan di sini hanya dirujuk ke
Bab 3 dan Bab 5.

Teorema Cantor--Bernstein--Schroeder di bab ini hanya sketsa. Pembuktian penuh
memerlukan pembukuan kedalaman rantai yang hati-hati; versi formal lengkapnya ada
pada Avigad--Massot, \emph{Mathematics in Lean}, bagian 4.3, sedangkan perlakuan
klasik untuk himpunan, fungsi, dan relasi ekuivalensi ada pada Hoffman dan Kunze,
\emph{Linear Algebra}, Apendiks A.1--A.3. Rujukan primer untuk diagonal dan
ketakterhitungan adalah Cantor, "Ueber eine Eigenschaft des Inbegriffes aller
reellen algebraischen Zahlen" (1874) dan "Ueber eine elementare Frage der
Mannigfaltigkeitslehre" (1891).

Pertanyaan terbuka yang paling terkenal: adakah kardinal di antara $|\mathbb{N}|$ dan
$|\mathbb{R}|$? Hipotesis kontinum menyatakan tidak ada. Godel (1940) menunjukkan
hipotesis itu konsisten dengan ZFC, Cohen (1963) menunjukkan ingkarannya juga konsisten.
Jadi pertanyaan itu tidak dapat dijawab di dalam ZFC.

Dua batas teknis. Pertama, $\sim$ tidak boleh disebut relasi ekuivalensi pada
"himpunan semua himpunan", sebab kumpulan itu bukan himpunan; perlakuan yang benar
memakai kelas sebagai objek metateori. Kedua, akibat "gabungan terhitung dari himpunan
terhitung adalah terhitung" memerlukan pemilihan terhitung bila penomorannya tidak
diberikan, dan Latihan 9 sengaja mensyaratkan penomoran eksplisit.
\end{untukpeneliti}

\peta{Bab 2 pada Abbott, \emph{Understanding Analysis}, bagian "Some Preliminaries"
(himpunan, fungsi, kardinalitas) dan "Cantor's Theorem"; Hoffman dan Kunze,
\emph{Linear Algebra}, Apendiks A.1 (himpunan), A.2 (fungsi dan invers), A.3 (relasi
ekuivalensi, kelas, kuosien); Avigad--Massot, \emph{Mathematics in Lean}, Bab 4
("Sets and Functions") untuk ekstensionalitas, prapeta, dan
Cantor--Bernstein--Schroeder, serta Bab 6.1 untuk himpunan berhingga.}
